\uuid{Df9j}
\titre{Théorème de Rolle}
\auteur{bodin ; hanani ; mzari}
\organisation{exo7}
\licence{CC BY-NC-SA 4.0}
\datecreate{2020-01-20}
\niveau{L1}
\module{Analyse}
\chapitre{Dérivabilité des fonctions réelles}
\sousChapitre{Autre}
\difficulte{3}
\theme{Dérivées, théorème de Rolle}
\exo7qcmid{lille-1.derivees.023}

\contenu{
\begin{question}
\qclassification{124.00}
\qtags{motcle=théorème de Rolle}
Soit $\displaystyle f(x)=x^3+x^2+x-\frac{13}{12}$. Quelles sont les bonnes réponses ?
\begin{answers}
    \bad{$\displaystyle f(0)=-\frac{13}{12}<0$ et $\displaystyle f(1)=-\frac{1}{12}<0$, donc $f(x)=0$  n'a pas de solution dans $]0,1[$.}
    \good{L'équation $f(x)=0$ admet une solution dans $]0,1[$.}
    \good{Le théorème de Rolle s'applique à une primitive de $f$ sur $[0,1]$.}
    \bad{Le théorème de Rolle s'applique à $f$ sur $[0,1]$.}
\end{answers}
\begin{explanations}
Le théorème de Rolle ne s'applique pas à $f$ sur $[0,1]$ car $\displaystyle f(0)\neq f(1)$. Mais on peut l'appliquer à
$$\displaystyle F(x)=\frac{x^4}{4}+\frac{x^3}{3}+\frac{x^2}{2}-\frac{13}{12}x.$$
Cette fonction vérifie toutes les hypothèses du théorème, donc
$$\exists c\in ]0,1[,\; F'(c)=0\Leftrightarrow f(c)=0.$$
\end{explanations}
\end{question}
}
